Diese Arbeit entwickelt und evaluiert einen erklaerbaren zweistufigen hierarchischen
\gb{}-Klassifikator zur Erkennung von Cyberangriffen in Internet of Things (\iot{})-Netzwerken
auf dem CICIoT2023-Datensatz~\cite{neto2023ciciot2023} (838.605~Instanzen).
Das Modell wird mit der \ba{} als primaerer Metrik bewertet; SHapley Additive Explanations
(\shap{}) dienen als post-hoc-Erklaerungskomponente und schliessen die bei baumbasierten
Modellen auf diesem Datensatz bestehende Erklaerbarkeitsluecke~\cite{raturi2026exploratory}.

Eine Ablation Study zeigt, dass die Reduktion mittels Principal Component Analysis (\pca{})
von 39 auf 16 Merkmale die \ba{} der Mehrklassenstufe um 8{,}21~Prozentpunkte senkt (0{,}5480
gegenueber 0{,}6301) und zugleich die semantische Interpretierbarkeit der
\shap{}-Erklaerungen verringert.
Eine klassengewichtete Trainingsstrategie hebt die \ba{} auf 0{,}7261; der Abstand zum von
Raturi et al.\ \cite{raturi2026exploratory} berichteten Referenzwert von 0{,}952 wird als
Reproduzierbarkeitsfrage eingeordnet und auf Unterschiede im Versuchsaufbau zurueckgefuehrt.
Die Qualitaet der \shap{}-Erklaerungen wird quantitativ (Faithfulness, Stabilitaet und
Jaccard-Aehnlichkeit nach Hermosilla et al.~\cite{hermosilla2025xai_forensic}) sowie qualitativ
ueber den Abgleich mit dokumentierten Angriffssignaturen bewertet.

\textbf{Schluesselwoerter:}
Internet of Things, Machine Learning, Gradient Boosting, CICIoT2023, Balanced Accuracy,
Intrusion Detection, Explainable Artificial Intelligence, SHAP, Principal Component
Analysis, Cyberangriffserkennung
